% !TEX program = xelatex
\documentclass[11pt,a4paper]{ctexart}

% ==================== 颜色 ====================
\usepackage{xcolor}
\definecolor{maincolor}{HTML}{1A7A6B}       % 主色（标题、顶栏、框边框）
\definecolor{accent}{HTML}{D4A843}          % 强调色（关键词高亮）
\definecolor{boxbg}{HTML}{F0F7F5}           % 问题框背景
\definecolor{borderblue}{HTML}{2E86AB}      % 方法框边框
\definecolor{proofbg}{HTML}{FDF8F0}         % 方法框背景
\definecolor{summarybg}{HTML}{F5F0FF}       % 总结框背景

% ==================== 页面 ====================
\usepackage[top=2.2cm, bottom=2.2cm, left=2.5cm, right=2.5cm, headheight=23.12pt]{geometry}

% ==================== 字体 ====================
\setmainfont{TeX Gyre Pagella}

% ==================== 数学 ====================
\usepackage{amsmath,amssymb,amsfonts,mathtools}

% ==================== 彩色框 ====================
\usepackage{tcolorbox}
\tcbuselibrary{skins,breakable,most}

% ==================== 页眉页脚 ====================
\usepackage{fancyhdr}
\pagestyle{fancy}
\fancyhf{}
\fancyhead[L]{\color{maincolor}\sffamily\bfseries 数学讲义}
\fancyhead[R]{\color{gray}\sffamily \today}
\fancyfoot[C]{\color{gray}\sffamily\thepage}
\renewcommand{\headrulewidth}{1.5pt}
\renewcommand{\headrule}{\hbox to\headwidth{\color{maincolor}\leaders\hrule height \headrulewidth\hfill}}

% ==================== 章节标题样式 ====================
\usepackage{titlesec}
\titleformat{\section}
  {\color{maincolor}\sffamily\LARGE\bfseries}
  {\thesection}{1em}{}
  [\vspace{0.3cm}{\color{maincolor!40}\titlerule[2pt]}]

\titleformat{\subsection}
  {\color{maincolor!80}\sffamily\Large\bfseries}
  {\thesubsection}{1em}{}

% ==================== 表格 ====================
\usepackage{array,booktabs,colortbl}

% ==================== 杂项 ====================
\usepackage{enumitem}
\setlist{nosep, leftmargin=*}
\usepackage[hidelinks]{hyperref}

% ==================== 自定义命令 ====================
\newcommand{\highlight}[1]{\colorbox{accent!20}{\sffamily\bfseries #1}}

% ==================== 自定义环境 ====================

% 1. 问题框 —— 青绿边框 + 浅绿背景
\newtcolorbox[auto counter]{problembox}[1][]{
  enhanced,
  colframe=maincolor,
  colback=boxbg,
  coltitle=white,
  fonttitle=\sffamily\bfseries,
  title={\large 问题},
  attach boxed title to top left={xshift=4mm, yshift=-2mm},
  boxed title style={colback=maincolor, sharp corners},
  sharp corners,
  boxrule=1.2pt,
  left=4mm, right=4mm, top=5mm, bottom=4mm,
  before skip=1.2em, after skip=1.2em,
  #1
}

% 2. 方法/证明框 —— 蓝边框 + 暖白背景
\newtcolorbox[auto counter]{methodbox}[2][]{
  enhanced,
  colframe=borderblue,
  colback=proofbg,
  coltitle=white,
  fonttitle=\sffamily\bfseries,
  title={#2},
  attach boxed title to top left={xshift=4mm, yshift=-2mm},
  boxed title style={colback=borderblue, sharp corners},
  sharp corners,
  boxrule=1.2pt,
  left=4mm, right=4mm, top=5mm, bottom=4mm,
  before skip=1.2em, after skip=1.2em,
  #1
}

% 3. 总结框 —— 淡紫背景
\newtcolorbox{summarybox}[1][]{
  enhanced,
  colframe=maincolor!60,
  colback=summarybg,
  coltitle=white,
  fonttitle=\sffamily\bfseries,
  title={\large $\blacktriangleright$ 总结},
  attach boxed title to top left={xshift=4mm, yshift=-2mm},
  boxed title style={colback=maincolor!70, sharp corners},
  sharp corners,
  boxrule=1.2pt,
  left=4mm, right=4mm, top=5mm, bottom=4mm,
  before skip=1.5em, after skip=1.2em,
  #1
}

% ==================== 文档开始 ====================
\begin{document}

% ---- 彩色顶栏标题 ----
\begin{tcolorbox}[
  enhanced,
  colframe=maincolor,
  colback=maincolor,
  sharp corners,
  boxrule=0pt,
  height=2.8cm,
  valign=center,
  before skip=-2.2cm,
  after skip=1.5cm,
  grow to left by=2.5cm,
  grow to right by=2.5cm,
]
  \begin{center}
    {\color{white}\sffamily\bfseries\Huge 求代数式最大值的两种方法}\\[3pt]
    {\color{white!85}\sffamily\large $\displaystyle \max \frac{xy + 2yz}{x^2 + y^2 + z^2}$}
  \end{center}
\end{tcolorbox}

\vspace{0.5cm}

\begin{flushright}
  {\color{gray}\sffamily \today}
\end{flushright}

% ==================== 正文 ====================

\section{问题陈述}

\begin{problembox}
求代数式 $\displaystyle \frac{xy + 2yz}{x^2 + y^2 + z^2}$ 的\highlight{最大值}。
\end{problembox}

本题可使用两种不同方法求解：方法一借助\highlight{基本不等式}与\highlight{待定系数法}，通过引入参数配凑系数；方法二则利用\highlight{柯西不等式}与\highlight{换元法}，从不同角度切入。

% ==================== 新页 ====================
\newpage

\section{解法一：基本不等式与待定系数法}

\begin{methodbox}{核心思路}
引入待定正数 $a,b>0$，对分子中的两项分别使用基本不等式 $uv \le \frac{\alpha}{2}u^2 + \frac{1}{2\alpha}v^2$ 进行放缩，再通过\highlight{令各项系数相等}的方式确定 $a,b$ 的值，从而配凑出分母并求得最大值。
\end{methodbox}

\subsection*{引入待定系数并放缩}

引入待定常数 $a>0,\; b>0$，对分子各项逐一使用基本不等式：

\[
xy \le \frac{a}{2}x^2 + \frac{1}{2a}y^2
\]

\[
2yz \le \frac{b}{2}y^2 + \frac{2}{b}z^2
\]

\subsection*{相加并配凑系数}

将上述两式相加，并按 $x^2,y^2,z^2$ 合并同类项：

\[
xy + 2yz \le \frac{a}{2}x^2 + \left(\frac{1}{2a} + \frac{b}{2}\right)y^2 + \frac{2}{b}z^2
\]

为配凑出分母 $x^2+y^2+z^2$ 的结构，令各项系数均等于同一常数 $k$：

\[
\frac{a}{2} = \frac{1}{2a} + \frac{b}{2} = \frac{2}{b} = k
\]

\subsection*{求解待定系数}

由连等式的左右两端可得：

\[
a = 2k,\quad b = \frac{2}{k}
\]

将其代入中间等式进行求解：

\[
\frac{1}{2(2k)} + \frac{2/k}{2} = k
\quad\Longrightarrow\quad
k^2 = \frac{5}{4}
\]

取正值 $k = \frac{\sqrt{5}}{2}$，即得：

\[
\boxed{\frac{xy + 2yz}{x^2 + y^2 + z^2} \le k = \frac{\sqrt{5}}{2}}
\]

因此该代数式的最大值为 $\displaystyle \frac{\sqrt{5}}{2}$。

% ==================== 新页 ====================
\newpage

\section{解法二：柯西不等式与换元法}

\begin{methodbox}{核心思路}
先将分子提取公因式 $y$，对含 $x,z$ 的部分使用\highlight{二维柯西不等式}，随后通过\highlight{换元} $t=\sqrt{x^2+z^2}$ 化简结构，最后再次使用\highlight{基本不等式}放缩分母即得结果。
\end{methodbox}

\subsection*{提取公因式与柯西不等式}

提取分子中的公因式 $y$：

\[
xy + 2yz = y(x + 2z)
\]

对含有 $x$ 与 $z$ 的部分使用二维柯西不等式：

\[
(x + 2z)^2 \le (1^2 + 2^2)(x^2 + z^2) = 5(x^2 + z^2)
\]

\subsection*{换元化简}

两边开平方，代回原分式中：

\[
\frac{y(x + 2z)}{x^2 + y^2 + z^2}
\le \frac{\sqrt{5}\,y\sqrt{x^2 + z^2}}{x^2 + y^2 + z^2}
\]

为简化结构，令 $t = \sqrt{x^2 + z^2}$ 进行整体换元，原式化为：

\[
\frac{\sqrt{5}\,yt}{y^2 + t^2}
\]

\subsection*{基本不等式放缩分母}

利用基本不等式 $y^2 + t^2 \ge 2yt$ 放缩分母：

\[
\frac{\sqrt{5}\,yt}{y^2 + t^2}
\le \frac{\sqrt{5}\,yt}{2yt}
= \frac{\sqrt{5}}{2}
\]

因此：

\[
\boxed{\frac{xy + 2yz}{x^2 + y^2 + z^2} \le \frac{\sqrt{5}}{2}}
\]

即该代数式的最大值为 $\displaystyle \frac{\sqrt{5}}{2}$。

\section{总结}

\begin{summarybox}

\subsection*{方法对比}

\begin{tabular}{@{}>{\sffamily\bfseries}p{2.8cm} p{5.5cm} p{5.5cm}@{}}
\toprule
& \color{maincolor} \textbf{解法一：基本不等式与待定系数法} & \color{borderblue} \textbf{解法二：柯西不等式与换元法} \\
\midrule
\textbf{切入点}   & 直接对分子各项引入参数放缩 & 提取公因式 $y$，分离 $x,z$ 结构 \\
\textbf{核心工具} & 基本不等式 $+$ 待定系数 & 柯西不等式 $+$ 换元 $+$ 基本不等式 \\
\textbf{关键技巧} & 令系数相等以配凑分母 & 令 $t=\sqrt{x^2+z^2}$ 化简分式 \\
\bottomrule
\end{tabular}

\vspace{1em}

\subsection*{结论}

两种解法从不同角度切入，最终求得该代数式的最大值均为：

\begin{center}
\begin{tcolorbox}[
  enhanced, colframe=accent!80, colback=accent!10,
  sharp corners, boxrule=1.5pt, width=6cm, center,
]
  {\color{maincolor}\sffamily\bfseries\LARGE $\displaystyle \frac{\sqrt{5}}{2}$}
\end{tcolorbox}
\end{center}

\end{summarybox}

\end{document}
